
\begin{remarque}
\label{V.2.5}
\ifthenelse{\boolean{orig}}
{La hipótesis de que~$G$ actúa fielmente no es evidentemente por superflua}
{La hipótesis de que~$G$ actúa fielmente evidentemente no es superflua}
en el Corolario~\Ref{V.2.4}. Lo mismo ocurre con la hipótesis de que~$X$ es
conexo, como se ve, por ejemplo, tomando $X=Y\times E$, donde~$E$ es un
conjunto finito, y tomando como~$G$ el grupo de permutaciones de~$E$: $G$
actúa con mucha inercia; no obstante,
$(Y\times E)/G=Y\times(E/G)=Y$, y~$X$ es étale sobre~$Y$. Tomando como~$G$
un grupo estrictamente menor que el grupo simétrico de~$E$, pero que actúe
transitivamente sobre~$E$, se ve que habrá también $Y$-automorfismos de~$X$
que no proceden de~$G$.
\end{remarque}

El ejemplo típico de un grupo~$G$ que actúa sin inercia es el de
$Y\times G$, sobre el cual se hace actuar~$G$ mediante sus traslaciones por
la derecha sobre el factor~$G$: un $Y$-preesquema~$X$ con grupo de operadores
por la derecha~$G$ se dice \emph{trivial} si es isomorfo a~$Y\times G$.

Para relacionar
\ifthenelse{\boolean{orig}}{el preesquemas}{los preesquemas}
con grupos finitos de operadores y la noción de fibrado principal en una
categoría (relación que no necesitaremos, por lo demás, en el resto del
seminario, pero que es importante en otros contextos), son útiles las
consideraciones siguientes. Fijamos un preesquema base~$Y$ y nos situamos
\ifthenelse{\boolean{orig}}{en el categoría}{en la categoría}
de los $Y$-preesquemas. Si~$G$ es un grupo
\marginpar{114}
finito, pondremos para abreviar $G_Y=Y\times G$; se trata, pues, de un
esquema en grupos finito sobre~$Y$ (\cf \numero 1), y si~$X$ es un
$Y$-preesquema, se tiene
$$
X\times_YG_Y=X\times G
$$
(misma referencia). Dar un $Y$-morfismo $X\times_YG_Y\to X$ equivale, por
tanto, a dar un $Y$-morfismo $X\times G\to X$, \ie a dar, para todo
$g\in G$, un $Y$-morfismo $T_g\colon X\to X$. Se comprueba inmediatamente
que, para que los~$T_g$ definan sobre~$X$ una estructura de preesquema con
grupo de operadores por la derecha~$G$ (\ie $T_{gg'}=T_{g'}T_g$ y
\ifthenelse{\boolean{orig}}{$T_e=\id_x$}{$T_e=\id_X$}),
es necesario y suficiente que el $Y$-morfismo correspondiente
$X\times_YG_Y\to X$ defina sobre~$X$ una estructura de $Y$-preesquema con
acción del $Y$-esquema en grupos~$G_Y$ (en el sentido general de los objetos
con acción de un $\cal{C}$-grupo en una categoría~$\cal{C}$). Supongamos que
así ocurre. Recordemos que~$X$ se dice \emph{formalmente homogéneo principal}
\index{formalmente homogéneo principal (preesquema)|hyperpage}%
bajo~$G_Y$\kern1pt\footnote{Ahora se dice más bien que~$X$ es un
pseudotorsor bajo~$G_Y$.} si el morfismo canónico
$$
X\times_YG_Y\to X\times_YX
$$
cuyas componentes son, respectivamente, $\pr_1$ y el morfismo de acción
$\pi\colon X\times_YG_Y\to X$, es un isomorfismo. En este caso,
identificando el primer miembro con~$X\times G$, el morfismo considerado es
el que asocia a todo $g\in G$ el morfismo
$$
(\id_X,T_g)=(\id_X\times_YT_g)\Delta_{X/Y}\colon X\to X\times_YX.
$$
Por consiguiente, decir que~$X$ es formalmente homogéneo principal bajo~$G_Y$
significa también que $X\times_YX$ es isomorfo a la suma disjunta de los
transformados de la diagonal por los elementos $(e,g)$ de~$G\times G$ (que
actúan sobre~$X\times_YX$ de manera evidente), donde~$e$ designa el elemento
neutro de~$G$. Si no se quiere distinguir entre izquierda y derecha y se
desea una fórmula que siga siendo aplicable al producto de más de dos
factores idénticos a~$X$, puede formularse la condición diciendo que el
morfismo canónico
$$
X\times_G(G\times G)\to X\times_YX
$$
obtenido asociando al par $(g,g')$ el morfismo
$$
(T_g,T_{g'})=(T_g\times_YT_{g'})\Delta_{X/Y}\colon X\to X\times_YX
$$
y haciendo actuar~$G$ por la izquierda sobre~$G\times G$ mediante el
homomorfismo diagonal
$$
s(g,g')=(sg,sg'),
$$
es un \emph{isomorfismo}.

La
\marginpar{115}
noción de \emph{espacio homogéneo principal}
\index{homogéneo principal (preesquema, fibrado)|hyperpage}%
se obtiene de la de espacio formalmente homogéneo principal añadiendo un
axioma suplementario que asegure que el «cociente» de~$X$ por~$G_Y$ existe y
es precisamente el objeto unidad a la derecha de la categoría, aquí~$Y$.
Este axioma puede variar según el contexto, y a menudo se explicita de la
manera más cómoda (en el yoga del «descenso») exigiendo que el objeto con
operadores se vuelva «trivial», \ie isomorfo al producto
$X\times_YG_Y$ (en este caso, $X\times G$), tras un cambio de base adecuado,
de tipo especificado (de modo que, en la práctica, permita la técnica del
descenso; \cf Grothendieck, \emph{Technique de descente et théorèmes
d'existence en Géométrie Algébrique}, Sém. Bourbaki \No190, pp. 26--28)%
\footnote{\Cf exp\ptbl\ \Ref{VIII} para la teoría del descenso plano.}.
En este orden de ideas, señalemos aquí la caracterización de los fibrados
homogéneos principales de grupo~$G$ (en el sentido de~\loccit):

\begin{proposition}
\label{V.2.6}
Sean~$Y$ un preesquema localmente noetheriano y~$X$ un $Y$-preesquema con un
grupo finito de operadores~$G$ que actúa por la derecha. Las condiciones
siguientes son equivalentes:
\begin{enumerate}
\item[(i)] $X$ es finito sobre~$Y$, $Y=X/G$, y los grupos de inercia de los
puntos de~$X$ se reducen al elemento neutro.
\item[(ii)] Existe un cambio de base fielmente plano y cuasicompacto
$Y_1\to Y$ tal que $X_1=X\times_YY_1$ sea un $Y_1$-preesquema con operadores
trivial, \ie isomorfo a~$Y_1\times G$.
\item[(ii~bis)] Como en (ii), pero siendo $Y_1\to Y$ finito, étale y
sobreyectivo.
\item[(iii)] $X$ es formalmente homogéneo principal bajo~$G_Y$, y es
fielmente plano y cuasicompacto sobre~$Y$.
\end{enumerate}
\end{proposition}
\textit{Demostración.} (i)$\To$(ii~bis). Se toma $Y_1=X$, observando que
$X\to Y$ es, en efecto, finito, étale por~\Ref{V.2.3} y sobreyectivo.
Demostremos que~$X_1$ es entonces trivial sobre~$Y_1$,
\ifthenelse{\boolean{orig}}
{lo que resultará del}
{como se deduce del}

\begin{corollaire}
\label{V.2.7}
Si se verifica \textup{(i)} y~$X$ admite una sección sobre~$Y$, entonces~$X$
es un espacio con operadores trivial.
\end{corollaire}

En efecto, esta sección permite definir un $G$-morfismo $\ifthenelse{\boolean{orig}}{X}{Y}\times G\to X$,
sobreyectivo porque~$G$ es transitivo sobre las fibras de~$X$, inyectivo
porque~$G$ actúa sin inercia y, por último, isomorfismo local en virtud
de~\Ref{I.5.3}, pues~$X$ es étale sobre~$Y$. Por tanto, es un isomorfismo.

(ii~bis)
\marginpar{116}
implica trivialmente (ii), que implica (i), pues los ingredientes de (i) son
«invariantes» por extensión fielmente plana y cuasicompacta de la base
(\cf el Seminario Bourbaki citado anteriormente para «finito»; para los
grupos de inercia se aplica~\Ref{V.2.1}, y para $Y=X/G$, un recíproco
de~\Ref{V.1.9} en el caso de un cambio de base \emph{fielmente plano}, que
habíamos olvidado explicitar).

\ifthenelse{\boolean{orig}}
{Hemos demostrado (i)$\To$(iii) al pasar demostrando (i)$\To$(ii~bis).}
{Al demostrar (i)$\To$(ii~bis), se demostró también
(i)$\To$(iii).}
Finalmente, (iii)$\To$(ii), pues la primera hipótesis de (iii) significa
precisamente que~$X$ se vuelve trivial al hacer el cambio de base $Y_1=X$;
de aquí (ii), ya que~$X$ es fielmente plano y cuasicompacto sobre~$Y$.

\begin{definition}
\label{V.2.8}
Un $Y$-preesquema~$X$ con grupo de operadores por la derecha~$G$ que satisface
las condiciones equivalentes de~\Ref{V.2.6} se llama \emph{recubrimiento
principal de~$Y$, de grupo de Galois~$G$}.
\index{grupo de Galois de un recubrimiento principal|hyperpage}%
\index{principal (recubrimiento)|hyperpage}%
\index{recubrimiento principal|hyperpage}%
\end{definition}
